\section{The extension to arbitrary finite fields}

Let $\mathbb F_q$ have characteristic $p$. The arguments of Section~2 extend
after replacing five by $p$. We state the results separately because the Lean
formalization currently treats only $\mathbb F_5$.

Fix nonzero polynomials $p_1,p_2,p_3,p_4\in\mathbb F_q[t]$ such that
$p_2(0)p_4(0)\ne0$, and impose $g(0)\ne0$, $f(0)=0$, and $f'(0)\ne0$. Set
\[
g(t)=\frac{p_1(t)}{p_2(t)},
\qquad
f(t)=\frac{p_3(t)}{p_4(t)}.
\]
The $k$th column generating function and its denominator before cancellation
are
\[
H_k(t)=\frac{p_1(t)p_3(t)^k}{p_2(t)p_4(t)^k},
\qquad
Q_k(t)=p_2(t)p_4(t)^k.
\]
Define
\[
\pi_k=\polyord(Q_k)
=\min\{N\geqslant1:Q_k(t)\mid t^N-1\}.
\]
The coefficients of $H_k$ are eventually $\pi_k$-periodic. If
$Q_k$ is coprime to $p_1p_3^k$, then $\pi_k$ is the least eventual period.

\subsection{Period growth}

The Frobenius identity is
\begin{equation}\label{eq:q-frobenius}
(t^N-1)^p=t^{pN}-1.
\end{equation}
As before, $Q_k\mid Q_{k+1}$ gives $\pi_k\mid\pi_{k+1}$. For
$k\geqslant1$, one also has
\[
Q_{k+1}\mid Q_k^2\mid Q_k^p.
\]
Together with Equation~\eqref{eq:q-frobenius}, this gives
$\pi_{k+1}\mid p\pi_k$.

\begin{prop}\label{prop:q-ratio}
If $k\geqslant1$, then
\[
\pi_{k+1}=\pi_k
\qquad\text{or}\qquad
\pi_{k+1}=p\pi_k.
\]
Consequently, there is a nondecreasing sequence $(u_k)_{k\geqslant1}$ with
$u_1=0$ such that
\[
\pi_k=\pi_1p^{u_k}.
\]
\end{prop}

For the explicit formula, define
\[
\lambda_p(m)=\min\{j\geqslant0:m\leqslant p^j\}
\qquad(m\geqslant1).
\]
Suppose a common finite family of pairwise nonassociated irreducibles gives
\[
p_2=\prod_{i\in I}\varphi_i^{a_i},
\qquad
p_4=\prod_{i\in I}\varphi_i^{b_i},
\]
where $\varphi_i(0)\ne0$ and $a_i+kb_i\geqslant1$. Write
$e_i=\polyord(\varphi_i)$. The roots of $\varphi_i$ lie in the multiplicative
group of a finite extension of $\mathbb F_q$, so $p\nmid e_i$. The
prime-power formula becomes
\[
\polyord(\varphi_i^m)=e_i p^{\lambda_p(m)}.
\]
Therefore
\begin{prop}\label{prop:q-formula}
Under the preceding factorization hypotheses,
\begin{align*}
\pi_k
&=\mathop{\operatorname{lcm}}_{i\in I}
  \left(e_i p^{\lambda_p(a_i+kb_i)}\right)\\
&=E p^{s_k},
\end{align*}
where
\[
E=\mathop{\operatorname{lcm}}_{i\in I}e_i,
\qquad
s_k=\max_{i\in I}\lambda_p(a_i+kb_i).
\]
\end{prop}

Thus rational second components produce a tower of changing block lengths.
This differs from the polynomial second-component setting, where the
denominator of the column generating function is fixed and one circulant
space suffices for all columns.

\subsection{The circulant correspondence}

For $N\geqslant1$, let
\[
\mathcal C_N(\mathbb F_q)=\mathbb F_q[t]/(t^N-1).
\]
Take $N=\pi_{k+1}$. Both $H_k$ and $H_{k+1}$ are eventually $N$-periodic
because $\pi_k\mid\pi_{k+1}$. If $B_k$ and $B_{k+1}$ are length-$N$ blocks
read from the same sufficiently late position, then
\[
p_4(t)H_{k+1}(t)=p_3(t)H_k(t)
\]
gives the following relation.

\begin{prop}\label{prop:q-block}
The neighboring periodic blocks satisfy
\[
p_4(t)B_{k+1}=p_3(t)B_k
\quad\text{in }\mathcal C_N(\mathbb F_q).
\]
The block $B_k$ is its length-$\pi_k$ block repeated $N/\pi_k$ times.
If $p_4$ is nonconstant, then its class in
$\mathcal C_N(\mathbb F_q)$ is not a unit.
\end{prop}

The last assertion follows from
\[
p_4\mid Q_{k+1}\mid t^N-1.
\]
Thus the block relation specifies the affine solution set
\[
\{B\in\mathcal C_N(\mathbb F_q):p_4B=p_3B_k\}
\]
rather than the value of an invertible circulant operator applied to $B_k$.

\subsection{Pascal's array modulo \texorpdfstring{$p$}{p}}

For
\[
g(t)=\frac{1}{1-t},
\qquad
f(t)=\frac{t}{1-t},
\]
the $k$th column has generating function
\[
H_k(t)=\frac{t^k}{(1-t)^{k+1}}.
\]
Since $\polyord((1-t)^m)=p^{\lambda_p(m)}$, its period is
\[
\pi_k=p^{\lambda_p(k+1)}.
\]
For $N=\pi_{k+1}$, sufficiently late blocks satisfy
\[
(1-t)B_{k+1}=tB_k
\quad\text{in }\mathbb F_q[t]/(t^N-1).
\]
The neighboring columns are therefore related by cyclic discrete
antidifferentiation. The failure of $1-t$ to be invertible accounts for the
nonuniqueness of a periodic antidifference when only the preceding block is
specified.
